\section{Flujo sugerido para una clase}

\subsection{Antes de que entren los alumnos}

\begin{enumerate}
  \item Levantar Docker y validar el health check.
  \item Confirmar que el listado de pacientes carga.
  \item Entrar como Docente desde el selector superior.
  \item Crear o revisar una regla de demostracion.
  \item Verificar que el sandbox del docente tiene actividad limpia.
\end{enumerate}

\subsection{Durante la clase}

\begin{enumerate}
  \item Compartir la URL del RCE.
  \item Pedir a cada alumno que abra \texttt{/patients}.
  \item Elegir un paciente adulto.
  \item Abrir su ficha clinica.
  \item Mostrar que no hay recomendaciones activas o revisar las existentes.
  \item Crear una regla CQL.
  \item Validar la regla y revisar el ELM.
  \item Probar la regla contra el paciente.
  \item Cambiar un dato clinico editable.
  \item Guardar y observar si aparece una card CDS.
  \item Abrir Actividad CDS para ver la trazabilidad del hook.
\end{enumerate}

\rceScreenshot{03-pacientes-listado.png}{Listado de pacientes sinteticos con filtros.}
\rceScreenshot{05-ficha-paciente.png}{Ficha clinica de un paciente sintetico.}
\rceScreenshot{06-editar-datos.png}{Drawer para editar datos controlados del paciente.}

\subsection{Que debe explicar el docente}

El punto importante no es que la card aparezca por magia, sino mostrar la cadena:

\begin{lstlisting}
Alumno edita CQL
  -> backend valida CQL contra CQL Translation Service
  -> backend guarda CQL/ELM como Library FHIR en HAPI
  -> alumno edita datos del paciente en su sandbox
  -> backend arma el bundle FHIR efectivo
  -> motor CQL evalua la regla
  -> CDS Hooks devuelve cards
\end{lstlisting}
